\chapter{Diskrete Informationsdynamik}

\section{Der diskrete Informationszustand}

Die IWT beschreibt jeden physikalischen Zustand durch eine diskrete Informationsverteilung. Diese wird durch eine skalare Dichtesequenz

\[
I_k^{(n)}
\]

repräsentiert, die angibt, wie viel strukturierte Information am Netzwerkknoten \( k \) zum Zeitschritt \( n \) vorliegt.

Dabei ist:

\begin{itemize}
    \item \( k \in \{1, 2, \dots, N\} \) – Index der diskreten Informationszelle (Knoten),
    \item \( n \in \mathbb{N}_0 \) – diskreter Zeitindex (Zeitschrift),
    \item \( I_k^{(n)} \in \mathbb{R}^+ \) – Informationswert (skalar, dimensionslos).
\end{itemize}

Im Gegensatz zu klassischen Feldern besitzt \( I_k^{(n)} \) keine materielle Bedeutung. Sie beschreibt weder Masse noch Ladung oder Energie, sondern die Organisation eines physikalischen Systems. Energie, Impuls und andere Größen entstehen erst als abgeleitete Funktionale dieser Informationsstruktur.

\section{Diskrete Informationserhaltung}

Analog zur Kontinuitätsgleichung der klassischen Physik wird der Informationsfluss durch diskrete Differenzengleichungen beschrieben. Die fundamentale Erhaltungsgleichung im diskreten Netz lautet:

\[
\sum_{k \in \mathcal{N}(l)} \left( I_k^{(n+1)} - I_k^{(n)} \right) + \sum_{m \in \partial \mathcal{N}(l)} J_{lm}^{(n)} = 0,
\]

wobei:

\begin{itemize}
    \item \( \mathcal{N}(l) \) – Menge der Nachbarknoten von \( l \),
    \item \( \partial \mathcal{N}(l) \) – Rand des Nachbarschaftsbereichs,
    \item \( J_{lm}^{(n)} \) – Informationsfluss von Knoten \( l \) zu \( m \) in Zeitschritt \( n \).
\end{itemize}

Diese Gleichung ist das diskrete Herzstück der Theorie: Sie ersetzt die kontinuierliche Energieerhaltung durch eine diskrete Informationserhaltung. Die gesamte Dynamik ergibt sich aus der rekursiven Umlagerung von Information zwischen Netzwerkknoten.

\section{Information als Ursprung physikalischer Größen}

In der diskreten IWT entstehen physikalische Größen als Funktionale der Informationsverteilung \( I_k^{(n)} \).

\subsection{Diskrete Energie}

Die Energie am Knoten \( k \) zum Zeitpunkt \( n \) ist:

\[
E_k^{(n)} = \alpha \left( I_k^{(n)} - I_k^{(n-1)} \right)^2 + \beta \sum_{l \in \mathcal{N}(k)} \left( I_k^{(n)} - I_l^{(n)} \right)^2,
\]

mit Kopplungskonstanten \( \alpha, \beta \).

\subsection{Diskreter Impuls}

Der Impulsfluss zwischen Knoten \( k \) und \( l \) ist:

\[
p_{kl}^{(n)} = \gamma \left( I_k^{(n)} - I_k^{(n-1)} \right) \left( I_l^{(n)} - I_l^{(n-1)} \right),
\]

mit Kopplungskonstante \( \gamma \).

\subsection{Diskrete Masse (Trägheit)}

Die effektive Masse am Knoten \( k \) ist:

\[
m_k^{(n)} = \delta \sum_{l \in \mathcal{N}(k)} \left( I_k^{(n)} - I_l^{(n)} \right)^2,
\]

mit Kopplungskonstante \( \delta \). Sie misst den Widerstand gegen Änderungen der lokalen Informationsstruktur.

Damit wird die klassische Unterscheidung zwischen Materie, Feldern und Geometrie aufgehoben: Alles entsteht aus einer einzigen fundamentalen diskreten Größe – der Information.

\section{Dynamik als diskreter Informationsfluss}

Die Bewegungsgleichungen eines Systems ergeben sich aus der rekursiven Umlagerung von Information. Die Theorie unterscheidet zwei komplementäre diskrete Dynamikformen:

\begin{itemize}
    \item \textbf{Lokale diskrete Dynamik:} beschrieben durch die diskrete Weber-Kraft,
    \item \textbf{Globale diskrete Dynamik:} beschrieben durch das diskrete Bohm-Potential.
\end{itemize}

Diese beiden Strukturen sind keine konkurrierenden Modelle, sondern zwei Projektionen derselben diskreten Informationsdynamik.

\subsection{Lokale diskrete Dynamik: Diskrete Weber-Kraft}

Die diskrete Weber-Kraft beschreibt lokale Informationsflüsse zwischen benachbarten Knoten. Für zwei wechselwirkende Knotengruppen \( A \) und \( B \) lautet sie:

\[
\vec{F}_{AB}^{(n)} = \frac{q_A q_B}{4\pi\varepsilon_0 (r_{AB}^{(n)})^2} \left[ 1 - \frac{1}{c^2} \left( \frac{r_{AB}^{(n)} - r_{AB}^{(n-1)}}{T} \right)^2 + \frac{2 r_{AB}^{(n)}}{c^2} \cdot \frac{r_{AB}^{(n+1)} - 2r_{AB}^{(n)} + r_{AB}^{(n-1)}}{T^2} \right] \hat{r}_{AB}^{(n)}.
\]

Eigenschaften:

\begin{enumerate}
    \item Explizit rekursiv: Berechnet \( r^{(n+1)} \) aus \( r^{(n)} \) und \( r^{(n-1)} \),
    \item Nicht-zirkulär: Keine implizite Abhängigkeit \( F = f(\vec{r}) \) mit \( \vec{r} = F/m \),
    \item Lokal: Nur Nachbarkopplungen,
    \item Feldlos: Benötigt keine kontinuierlichen Felder.
\end{enumerate}

\subsection{Globale diskrete Dynamik: Diskrete Bohm-Struktur}

Das diskrete Bohm-Potential beschreibt die systemische, nicht-lokale Organisation des Informationszustands. Im diskreten Netz lautet es:

\[
Q_k^{(n)} = -\frac{\hbar^2}{2m} \frac{\Delta_d^2 \sqrt{I_k^{(n)}}}{\sqrt{I_k^{(n)}}},
\]

mit dem diskreten Laplace-Operator:

\[
\Delta_d^2 f_k = \sum_{l \in \mathcal{N}(k)} (f_l - f_k).
\]

Eigenschaften:

\begin{enumerate}
    \item Global: Wirkt über das gesamte Netz,
    \item Instantan: Keine Retardierung,
    \item Nicht-energetisch: Transportiert keine Energie,
    \item Organisierend: Optimiert die Gesamtstruktur.
\end{enumerate}

\subsection{Die digitale WDBT als fundamentales Netzwerk}

Die digitale WDBT beschreibt die Gesamtwirkung auf einen Knoten \( k \) durch drei diskrete Beiträge:

\[
F_k^{(n)} = F_{\text{WED},k}^{(n)} + F_{\text{WG},k}^{(n)} + F_{Q,k}^{(n)},
\]

wobei:

\begin{itemize}
    \item \( F_{\text{WED},k}^{(n)} \) – diskrete Weber-Elektrodynamik (Ladungen),
    \item \( F_{\text{WG},k}^{(n)} \) – diskrete Weber-Gravitation (Massen),
    \item \( F_{Q,k}^{(n)} \) – diskrete Bohm-Struktur (globale Organisation).
\end{itemize}

Diese digitale Theorie besitzt kein vorgegebenes Raummodell. Der Raum emergiert aus der Kopplungsstruktur \( K_{kl}^{(n)} \).

\section{Zusammenfassung}

Die diskrete Informationsdynamik der IWT basiert auf:

\begin{itemize}
    \item einem diskreten Informationszustand \( I_k^{(n)} \) als Grundgröße,
    \item einer diskreten Erhaltungsgleichung für Information,
    \item einer zweistufigen Dynamik: lokal (Weber-Kraft) und global (Bohm-Potential),
    \item der Emergenz von Raum, Zeit, Energie, Impuls und Masse aus der Informationsstruktur.
\end{itemize}

Damit ist die fundamentale Dynamik der IWT vollständig definiert, ohne Rückgriff auf kontinuierliche Felder oder eine vorgegebene Raumzeit.